\chapter*{Введение}
\addcontentsline{toc}{chapter}{Введение}
\label{cha:intro}
%\Introduction

Паллетизация остаётся ключевым технологическим узлом современной складской логистики: именно от плотности и устойчивости укладки зависит стоимость хранения, расход упаковочных материалов и объём перевозимого груза. С ростом ассортимента e-commerce-складов задача становится всё более вариативной, а ручная разработка правил паллетизации — трудоёмкой и часто приводит к недогрузам до 50 \%.

\textbf{Цель работы} --- разработка и программная реализация семейства алгоритмов трёхмерной укладки коробок на паллету, обеспечивающих максимальный коэффициент заполнения при гарантированной устойчивости размещения в ограниченное время.

В первом разделе работы содержится постановка задачи и подзадач, на которые она была разбита.

Во втором разделе работы приведён обзор научно-технической литературы по теме исследований.

В третьем разделе формализуется задача и описываются исходные данные, используемые для построения алгоритмов.
 
В четвёртом разделе приведены описания разработанных алгоритмов: жадного алгоритма GreedySolver, эвристик DBLF и FFD, метаэвристического алгоритма LNSSolver на основе поиска в большой окрестности и генетического алгоритма GeneticSolver. Описан мультипаллетный режим с метрикой баланса высот.

В пятом разделе приведён сравнительный анализ разработанных алгоритмов на 436 реальных задачах в однопаллетном и мультипаллетном режимах.